% =======================================================
% 4. DEFINITION OF DONE
% =======================================================
\section{Definition of Done (DoD)}
\label{sec:dod}

La Definition of Done è la lista di criteri che un task deve soddisfare per essere
considerato completato.
Si applica a due livelli distinti e si \textbf{espande} al crescere del progetto: la
Sez.~\ref{sec:dod-doc} è attiva fin dal primo sprint, mentre, la Sez.~\ref{sec:dod-codice} dall'inizio dello sviluppo del PoC\textsubscript{G}.

% ── 4.1 Livello task — fase documentale ─────────────────────────────────────
\subsection{Livello task -- Fase documentale}
\label{sec:dod-doc}
\textit{Attiva da: Sprint 1. Applicabile a: tutti i task documentali.}

\subsubsection*{Creazione e tracciabilità}
\begin{itemize}
    \item Il task è stato creato tramite Google Form dal Responsabile durante la
          riunione del martedì.
    \item Il task è linkato nel verbale della riunione dal Redattore\textsubscript{G}.
    \item È stato creato un branch\textsubscript{G} da \texttt{development} direttamente dalla issue\textsubscript{G},
          con naming: \\
          \texttt{id-gestionale|tecnico\_esterno|interno\_nomeDoc}
    \item Il branch è stato pushato entro la fine del giorno in cui si inizia a
          lavorare (anche se il documento è vuoto), così la issue risulta
          \textit{In Progress} nel board.
\end{itemize}

\subsubsection*{Redazione}
\begin{itemize}
    \item Il documento rispetta il template\textsubscript{G} e la struttura concordata nelle Norme di
          Progetto.
    \item Il contenuto è completo rispetto allo scope definito nel task.
    \item Nessun segnaposto \texttt{[da compilare]} lasciato senza motivazione
          esplicita.
    \item Ortografia e grammatica verificate.
    \item Terminologia coerente con il Glossario di progetto.
\end{itemize}

\subsubsection*{Commit\textsubscript{G} e Pull Request\textsubscript{G}}
\begin{itemize}
    \item Ogni commit fa riferimento alla issue con \texttt{\#id} nel messaggio.
    \item I messaggi di commit sono concisi e fanno riferimento al documento di
          riunione se pertinente.
    \item La PR è aperta \textbf{entro domenica} per dare tempo al Verificatore\textsubscript{G} di
          revisionare prima del martedì successivo.
    \item La PR linka la issue corrispondente.
    \item La PR descrive brevemente le modifiche apportate.
\end{itemize}

\subsubsection*{Revisione}
\begin{itemize}
    \item Almeno un membro del team (il Verificatore assegnato) ha revisionato il
          documento.
    \item Tutti i commenti della review sono stati risolti o discussi prima del merge\textsubscript{G}.
    \item È \textbf{il Redattore} a fare il merge, non il Verificatore.
\end{itemize}

\subsubsection*{Chiusura}
\begin{itemize}
    \item Merge effettuato su \texttt{main\textsubscript{G}}.
    \item Issue chiusa nel board.
    \item Ore effettive aggiornate su Google Form.
\end{itemize}

% ── 4.2 Livello task — fase sviluppo/PoC ────────────────────────────────────
\subsection{Livello task -- Fase sviluppo / PoC}
\label{sec:dod-codice}
\textit{Attiva da: Sprint N [da definire all'avvio dello sviluppo].
Applicabile a: task di sviluppo software.}

\medskip
\textbf{[Da compilare]}
\textit{Questa sezione verrà completata quando il team avvierà lo sviluppo del
Proof of Concept\textsubscript{G}. I criteri includeranno almeno:}
\begin{itemize}
    \item codice funzionante in locale senza errori;
    \item test unitari e/o di integrazione scritti e superati;
    \item copertura dei test non inferiore alla soglia concordata (es.\ $\geq 70\%$);
    \item nessun bug bloccante aperto;
    \item linting superato secondo la configurazione del repository\textsubscript{G};
    \item PR revisionata e approvata da almeno un membro del team;
    \item documentazione tecnica (commenti, Swagger/JSDoc) aggiornata;
    \item task chiuso nel board e ore rendicontate.
\end{itemize}
\textit{I criteri definitivi verranno concordati dal team prima dell'avvio dello
sprint di sviluppo e aggiunti qui.}

% ── 4.3 Livello sprint ───────────────────────────────────────────────────────
\subsection{Livello sprint}
\label{sec:dod-sprint}
Uno sprint è \textbf{Done} se:
\begin{itemize}
    \item tutti i task pianificati sono Done (Livello task), oppure esplicitamente
          rimandati al backlog con motivazione scritta nel verbale;
    \item il verbale della riunione del martedì è redatto e approvato;
    \item il registro dei rischi è stato rivalutato nella Sprint Review;
    \item la Retrospettiva ha prodotto almeno un'azione correttiva concreta per il
          prossimo sprint;
    \item le ore effettive di tutti i membri sono rendicontate su Google Form.
\end{itemize}

% ── 4.4 Checklist PR ────────────────────────────────────────────────────────
\subsection{Checklist PR -- Template GitHub}
\label{sec:dod-pr}
Il file \texttt{.github/pull\_request\_template.md} nel repository contiene la
seguente checklist, inserita automaticamente in ogni Pull Request:

\begin{verbatim}
Checklist DoD
- [ ] Documento completo e conforme al template (Norme di Progetto)
- [ ] Nessun [da compilare] senza motivazione esplicita
- [ ] Terminologia coerente con il Glossario
- [ ] Ogni commit ha #id nel messaggio
- [ ] PR aperta entro domenica
- [ ] PR linka la issue
- [ ] Revisione effettuata dal Verificatore assegnato
- [ ] Tutti i commenti della review risolti
- [ ] Merge fatto dal Redattore (non dal Revisore)
- [ ] Issue chiusa nel board
- [ ] Ore effettive rendicontate su Google Form
\end{verbatim}